\documentclass[10pt,letterpaper]{article}
\usepackage[utf8]{inputenc}
\usepackage[english]{babel}
\usepackage[margin=0.75in]{geometry}
\usepackage{titlesec}
\usepackage{enumitem}
\usepackage{hyperref}
\usepackage{xcolor}
\usepackage{mathptmx}

\definecolor{primaryColor}{RGB}{20, 50, 90}
\definecolor{secondaryColor}{RGB}{70, 80, 95}
\definecolor{lineColor}{RGB}{180, 190, 200}

\hypersetup{
    colorlinks=true,
    linkcolor=primaryColor,
    urlcolor=primaryColor,
    pdftitle={Curriculum Vitae - Rosh Guadiana Paez},
    pdfauthor={Rosh Guadiana Paez}
}

\titleformat{\section}{\large\bfseries\color{primaryColor}}{}{0em}{}[\color{lineColor}\titlerule]
\titlespacing*{\section}{0pt}{12pt}{6pt}
\setlist[itemize]{leftmargin=1.5em, topsep=2pt, itemsep=3pt, parsep=0pt, partopsep=0pt}

\begin{document}
\pagestyle{plain}

\begin{center}
    {\LARGE \bfseries \color{primaryColor} Rosh Guadiana P\'aez} \\[6pt]
    \color{secondaryColor}
    Puebla, Mexico $\cdot$ \href{mailto:rosh.guadianapz@udlap.mx}{rosh.guadianapz@udlap.mx} $\cdot$ +52 (618) 292-5697 \\
    \href{https://github.com/RoshTzsche}{github.com/RoshTzsche} $\cdot$ \href{https://linkedin.com/in/rosh-guadiana}{linkedin.com/in/rosh-guadiana}
\end{center}

\vspace{10pt}

\section{About Me}
Biomedical Engineering student with a robust analytical foundation in pure mathematics, computational modeling, and nonlinear dynamical systems. My expertise lies in biosignal processing, full-stack application development, and the integration of artificial intelligence in healthcare. I am highly focused on solving complex problems in neuroengineering.

\section{Academic Background}
\textbf{Universidad de las Am\'ericas Puebla (UDLAP)} \hfill San Andr\'es Cholula, Puebla \\
\textit{Bachelor of Science in Biomedical Engineering} \hfill Expected Graduation: 2028
\begin{itemize}
    \item \textbf{GPA:} 9.7/10 | \textbf{Funding:} Awarded a 100\% Academic Scholarship for sustained academic excellence.
    \item \textbf{STEM Representative Team:} Active member of the official UDLAP STEM team, contributing to institutional technological projects and science communication.
    \item \textbf{Advanced Training:} Selected institutional representative and participant in the \textit{Putnam Mathematical Competition} preparatory program (December 2026).
    \item \textbf{Relevant Coursework:} Ordinary Differential Equations, Chemical Thermodynamics, Calculus II, Electronics.
\end{itemize}

\section{Neuroengineering \& Technical Projects}

\textbf{Phase Space Attractor Reconstruction (EEG)} $|$ \href{https://github.com/RoshTzsche/signals_attractor}{Dashboard} \hfill \textit{Computational Neuroscience}
\begin{itemize}
    \item Analyzed Electroencephalography (EEG) time-series data using nonlinear dynamical systems theory.
    \item Reconstructed phase space attractors utilizing Takens' embedding theorem to map the geometric structure and chaotic dynamics of neurological activity.
\end{itemize}

\textbf{Hackathon por Amor a Puebla} \hfill \textit{Full-Stack \& Accessibility Developer} \hfill 2026
\begin{itemize}
    \item Developed \textit{Pipo} (\href{https://github.com/RoshTzsche/Pipo}{github.com/RoshTzsche/Pipo}), an automated medical classification application, handling the complete end-to-end technical implementation.
    \item Built backend architecture with extensive API integrations, programmed the interactive frontend interface utilizing modern web frameworks, and ensured strict accessibility compliance for healthcare users.
\end{itemize}

\textbf{Embedded Signal Processing \& PPG Sensor Design} \hfill \textit{Hardware Architecture}
\begin{itemize}
    \item Designed and implemented analog circuitry for a Photoplethysmography (PPG) sensor to acquire cardiovascular pulse data.
    \item Developed codebases for PIC16F1717 microcontrollers to execute digital filtering and signal conditioning, optimizing the signal-to-noise ratio of the acquired physiological data.
\end{itemize}

\textbf{NASA Space Apps Challenge} \hfill \textit{Developer \& Data Analyst} \hfill Two Consecutive Years
\begin{itemize}
    \item \textbf{Asteroid-DB (Recent Edition):} Developed a platform (\href{https://github.com/RoshTzsche/Asteroid-DB}{github.com/RoshTzsche/Asteroid-DB}) integrating and analyzing data from Near-Earth Object (NEO) programs to model astrophysical trajectories.
    \item \textbf{ExoSky (Previous Edition):} Modeled and analyzed exoplanet data and celestial mechanics.
\end{itemize}

\section{Research Experience}

\textbf{Undergraduate Researcher in Metagenomics} \hfill \textit{Universidad de las Am\'ericas Puebla} \\
\textit{Project: Dietary Metagenomics of the Narrow-Clawed Crayfish} \hfill 2026 -- Present
\begin{itemize}
    \item Implement optimized bioinformatics pipelines in Python and R for taxonomic classification and multidimensional diversity modeling (PCoA) under the supervision of Dr. Osiris D\'iaz Torres.
\end{itemize}
\newpage 
\section{Teaching, Outreach \& Science Communication}

\textbf{STEMpire Podcast (UDLAP STEM Team)} \hfill \textit{Interviewer, Editor \& Guest Speaker} \\
\textit{Science Communicator} \hfill Multiple Appearances
\begin{itemize}
    \item Core contributor to the official podcast of the UDLAP STEM Representative Team, editing episodes and engaging in scientific dissemination regarding biomedical engineering and research methodologies.
\end{itemize}

\textbf{Verano STEM} \hfill \textit{Universidad de las Am\'ericas Puebla (UDLAP)} \\
\textit{Guest Lecturer \& Instructor} \hfill 2025 -- Present
\begin{itemize}
    \item \textbf{Upcoming (Summer 2026):} Scheduled to deliver specialized lectures on Quantum Computation and Pure Mathematics for the institutional STEM outreach program.
    \item \textbf{Previous Editions:} Conducted instructional sessions and lectures focusing on the structural biology and fundamentals of DNA.
\end{itemize}

\textbf{Noche de las Estrellas} \hfill \textit{Workshop Instructor} \hfill 2024 \& 2025
\begin{itemize}
    \item Designed and facilitated interactive science workshops for massive public events aimed at astronomical and scientific dissemination.
\end{itemize}

\section{Conferences, Olympiads \& Additional Experience}
\begin{itemize}
    \item \textbf{Mathematical Competitions:} Active participant in the \textit{Olimpiada Mexicana Universitaria de Matem\'aticas} (OMUM 2026) and upcoming participant in the \textit{Pierre Fermat} competition (2026).
    \item \textbf{MUTVI 4:} Presented original computational work on the reconstruction of phase space attractors from Electroencephalography (EEG) signals.
    \item \textbf{Swift Social Change Makers Enactus (2025):} 2nd Place National Award for designing high-impact AI technological solutions.
    \item \textbf{National Physics Olympiad (2023):} Selected for the Durango State Delegation to compete in the National phase.
    \item \textbf{XXXII National Chemistry Olympiad (2022):} Awarded the National Bronze Medal.
    \item \textbf{State Biology Olympiad (2020):} Awarded the State Gold Medal.
    \item \textbf{COCYTED Advanced Training (2017 -- 2021):} Selected by the State Council of Science and Technology of Durango (COCYTED) for intensive extracurricular training in biology, chemistry, and medicine.
\end{itemize}

\section{Skills and Competencies}
\begin{itemize}
    \item \textbf{Programming \& Web Stack:} Python, C, R, Swift, React 18, Node.js, Bash.
    \item \textbf{Advanced Mathematics (Independent Study):} Real analysis, general topology, measure theory ($\mathcal{L}^p$ spaces), formal logic, and nonlinear dynamical systems theory.
    \item \textbf{Engineering \& Bioinformatics:} EEG data analysis, analog circuit design, digital signal processing, Kraken pipelines.
    \item \textbf{Working Environment:} Arch Linux ecosystem (CachyOS, Hyprland), Neovim, Git, \LaTeX\ / Typst.
    \item \textbf{Languages:} Spanish (Native), English (Advanced), French (Basic).
\end{itemize}

\end{document}
